\section[Introduction]{Introduction}

\begin{frame}{Problématique}
	\begin{itemize}
		\item Jusqu'à présent nous avons vu comment les modèles peuvent faire une prédiction pour des entrées $X$. Et globalement comment les modèles se comportent grâce à différentes métriques. Ou même comment localement on peut décomposer l'erreur sur des données de test. 
		\pause
		\item Mais pour une prédiction du modèle, dont nous ne connaissons pas le résultat vrai, quelle est l'erreur commise ?
		\pause
		\item Dans tous les domaines d'application, une erreur de prédiction a un \textbf{coût}.
	\end{itemize}
\end{frame}

\begin{frame}{Problématique}
	\begin{block}{Problématique}
		Tous les modèles sont des approximations. Essentiellement, tous les modèles sont faux, mais certains sont utiles. Cependant, il ne faut jamais perdre de vue le caractère approximatif du modèle. \\
		\vspace{0.5cm}
		\hfill
		\textit{--- Georges E. P. Box (2007)}
	\end{block}
	\vfill 
	\tiny \textcolor{gray}{Source: Response Surfaces, Mixtures, and Ridge Analyses - George E. P. Box, Norman R. Draper  - 2007 - p.414, John Wiley \& Sons}
\end{frame}

\begin{frame}{Régions de prédiction et intervalles de confiance}
	Dans ce cours nous introduirons les notions de régions de prédiction et d'intervalles de confiance, à la fois pour des problèmes de classification et pour des problèmes de régression.
	\begin{itemize}
		\item Dans un premier temps on développera l'intuition sur ces intervalles avant d'en donner une définition.
		\pause
		\item Ensuite nous verrons que les modèles linéaires (régression et classification) proposent des estimations théoriques de ces régions et intervalles.
		\pause
		\item Enfin, nous verrons comment, empiriquement, on peut établir ces régions et intervalles pour une grande variété de modèles.
	\end{itemize}
\end{frame}